\documentclass[10pt,a4paper]{amsart}
\usepackage[margin=19mm]{geometry}
\usepackage{amsmath,amssymb,mathtools}
\usepackage[scr=rsfs,cal=euler]{mathalfa}
\usepackage{xfrac}
\usepackage[unicode,hidelinks,bookmarks=false]{hyperref}
\newcommand{\ie}{i.e.\ }
\newcommand{\cf}{cf.\ }
\newcommand{\resp}{respectively\ }
\newcommand{\Subsection}{\subsection}
\providecommand{\nobreakdash}{\nobreak\hbox{-}}
\setlength{\emergencystretch}{2em}
\multlinegap0pt
\usepackage{amssymb}
\usepackage{xcolor}
\usepackage{algorithm}\usepackage{algpseudocode}
\newtheorem{theorem}{Theorem}
\title{Verification Framework for the Union of Control Barrier Functions}
\author{Chuanrui Jiang, Andrew Clark}
\date{Source: arXiv:2606.07892v1 (CC BY 4.0)}
\begin{document}
\maketitle

\noindent\emph{Source and licence.} Verification Framework for the Union of Control Barrier Functions. Version arXiv:2606.07892v1 (CC BY 4.0), distributed by arXiv under the Creative Commons Attribution 4.0 International licence, http://creativecommons.org/licenses/by/4.0/, as recorded in the arXiv licence metadata for this exact version and shown on the abstract page https://arxiv.org/abs/2606.07892v1. Full licence notice: \texttt{licenses/arxiv-2606.07892-CC-BY-4.0.txt}.\par\smallskip

\noindent\textit{Editorial scope.} Counted unit: the article's theorem 1 (source0.tex lines 268--276). The article's complete original proof of it is delivered with it (source0.tex lines 278--303, one \texttt{proof} environment of 26 lines). No mathematics is written, repaired or re-derived: the statement and the proof are the article's own lines, in the article's own notation, and no retained line is substituted or re-flowed.  Retained uncounted as the source-local context the proof needs: lines 108--111; lines 236--242.  Nothing was omitted from any retained span, so the package carries no omission record.  No retained line contains a citation, so the wrapper carries no bibliography and no bibliography entry.  No retained line contains a figure environment, an image-inclusion command or drawing code, so no image is delivered and the package declares no asset dependency.  Editorial formatting only, in addition: an AMS \texttt{amsart} preamble that restates the article's own declarations and macros used by the retained lines, namely the environments \texttt{theorem} on separate counters, together with the standard packages those lines need.  The retained environments print in this document as Theorem 1 in order of appearance, whereas the article prints the same items with the numbering of its own full document.  The complete native source of the article as distributed in the arXiv e-print for this exact version is bundled unchanged as \texttt{sources/202288/source0.tex}.
\begin{equation}
\label{eq:sys_dyn_control_aff}
    \Dot{\mathbf{x}}(t) = \mathbf{f}(\mathbf{x}(t)) + \mathbf{G}(\mathbf{x}(t))\mathbf{u}(t)
\end{equation}

\begin{theorem}
\label{thm:positive_switching_dist}
     Let $f: \mathbb{R}^{n_x}\rightarrow\mathbb{R}$ be a locally Lipschitz function with Lipschitz constant $L$ on a compact set $\mathcal{W}\subseteq\mathbb{R}^{n_x}$. Let $\eta > 0$ be a constant. For all $\mathbf{x}, \mathbf{x}'\in\mathcal{W}$, we have that 
     \begin{equation*}
         \inf \left\{ ||\mathbf{x}'-\mathbf{x}|| : f(\mathbf{x})\geq \eta, f(\mathbf{x}')\leq 0 \right\} > \frac{\eta}{L}
     \end{equation*}
\end{theorem}

\begin{theorem}
\label{thm:union-CBF-switching 1}
     Let $h_1(\mathbf{x})$, $\dots$, $h_{N_h}(\mathbf{x})$ be given differentiable functions with their first order derivatives being locally Lipschitz on $\mathcal{X}_c$. {Also let $\epsilon_{cbf}>0$ be a small constant and $\epsilon_{cbf}\leq \eta_l$.} If $\forall \mathbf{x}\in\mathcal{X}_c$, $\exists p\in\mathcal{P}, \mathbf{u}\in Int(\mathcal{U})$ with 
     \begin{itemize}
         \item[\textbf{C1-1}] $h_p(\mathbf{x})\geq 0$.
         \item[\textbf{C1-2}] $L_\mathbf{f}h_p(\mathbf{x}) + L_{\mathbf{G}}h_p(\mathbf{x})\mathbf{u}+\kappa h_p(\mathbf{x}) \geq \epsilon_{cbf}$.
     \end{itemize}
    Then, any switching signal $\sigma(t)$ following conditions S1-1 to S1-3 satisfies P1-1 and P1-2.
\end{theorem}

\begin{proof}
    We first prove P1-1 and then prove P1-2.

    \textit{\underline{Proof of P1-1: }} Assume $\sigma(t)$ switches to $p$ at time $t_k$, and the next switch of $\sigma(t)$ from $p$ to another index in $\mathcal{P}$ happens at $t_{k+1}$. In order to prove P1-1, we first prove that there exists a uniform lower bound for $||\mathbf{x}(t_{k+1}) - \mathbf{x}(t_k)||$, then we show that there exists a uniform lower bound on $t_{k+1}-t_{k}$ that holds for all $k=0,1,\ldots$. According to the problem setup, $\mathbf{f}(\mathbf{x})$ and $\mathbf{G}(\mathbf{x})$ of system \eqref{eq:sys_dyn_control_aff} are locally Lipschitz for all $\mathbf{x}\in\mathcal{X}$. 
    For all $p\in\mathcal{P}$, we let $f_p(\mathbf{x})=L_\mathbf{f}h_p(\mathbf{x}) + L_{\mathbf{G}}h_p(\mathbf{x})\mathbf{u}_p+\kappa h_p(\mathbf{x})$, where $\mathbf{u}_p\in\mathcal{U}$ is an admissible control such that $f_p(\mathbf{x}(t), \mathbf{u}_p)\geq \eta_h$. Since $\mathcal{U}$ is a compact set, we have $||\mathbf{u}_p||\leq U$ where $U=\max_{\mathbf{u}\in\mathcal{U}}||\mathbf{u}||$. We denote $L_{p}^{(G)}$ and $L_{p}^{(f)}$ be the Lipschitz constants of terms $(L_\mathbf{f}h_p(\mathbf{x})+\kappa h_p(\mathbf{x}))$ and $L_{\mathbf{G}}h_p(\mathbf{x})$ in the compact region $\mathcal{R}$, respectively. Hence, for all $\mathbf{x}\neq\mathbf{x}'$, we have:
    \begin{align*}
         \|f_p(\mathbf{x}) - f_p(\mathbf{x}')\|
         \leq (L_p^{(G)} + L_p^{(f)}U)\|\mathbf{x}-\mathbf{x}'\|
    \end{align*}
    We let $L_p = L_p^{(G)} + L_p^{(f)}U$ and $L=\max_{p\in\mathcal{P}}L_p$. 
    % According to condition S1-3 of Switching Strategy 1, at $\mathbf{x}(t_k)$, there exists a $\mathbf{u}_0\in Int(\mathcal{U})$ such that $L_\mathbf{f}h_p(\mathbf{x}(t_k)) + L_{\mathbf{G}}h_p(\mathbf{x}(t_k))\mathbf{u}_0+\kappa h_p(\mathbf{x}(t_k)) \geq \eta_h$. We now define $f_{p0}(\mathbf{x}) := L_\mathbf{f}h_p(\mathbf{x}) + L_{\mathbf{G}}h_p(\mathbf{x})\mathbf{u}_0+\kappa h_p(\mathbf{x})$. Since $\mathbf{f}(\mathbf{x})$, $\mathbf{G}(\mathbf{x})$ and $\partial h_p(\mathbf{x})/\partial\mathbf{x}$ are locally Lipchitz on $\mathcal{X}$, then the function $f_{p0}(\mathbf{x})$ is also locally Lipschitz on $\mathcal{X}$. {We let $L_0$ be the largest Lipschitz constant of $f_{p0}(\mathbf{x})$ among all $p\in\mathcal{P}$}. 
    By Theorem \ref{thm:positive_switching_dist}, for any $\mathbf{x}'\in\mathcal{X}_p$ such that $f_{p}(\mathbf{x}')\leq \eta_l$, we have $||\mathbf{x}' - \mathbf{x}(t_k)||\geq \frac{\eta_h-\eta_l}{L}$. Define the following two sets: 
    \begin{align*}
        \mathcal{Y} &:= \{\mathbf{x}|\sup_{\mathbf{u}\in Int(\mathcal{U})}\{L_\mathbf{f}h_p(\mathbf{x}) + L_{\mathbf{G}}h_p(\mathbf{x})\mathbf{u}+\kappa h_p(\mathbf{x})\}\leq \eta_l\}\\
        \mathcal{Y}_0 &:= \{\mathbf{x}| L_\mathbf{f}h_p(\mathbf{x}) + L_{\mathbf{G}}h_p(\mathbf{x})\mathbf{u}_p+\kappa h_p(\mathbf{x}) \leq \eta_l\}
    \end{align*}
    where $\mathbf{u}_0\in Int(\mathcal{U})$. It is obvious that $\mathcal{Y}\subseteq\mathcal{Y}_0$. Since we have $\mathbf{x}(t_k)\notin\mathcal{Y}$, $\mathbf{x}(t_k)\notin\mathcal{Y}_0$, and $\mathbf{x}(t_{k+1})\in\mathcal{Y}$, we have that
    \begin{equation*}
        ||\mathbf{x}(t_{k+1})- \mathbf{x}(t_k)||\geq||\mathbf{x}'- \mathbf{x}(t_k)||\geq {\frac{\eta_h-\eta_l}{L}}
    \end{equation*}
    Hence, we have a uniform lower bound for the term $||\mathbf{x}(t_{k+1})- \mathbf{x}(t_k)||\geq \frac{\eta_h-\eta_l}{L}$. {Since we also assume that $\mathcal{R}$ and $\mathcal{U}$ are bounded sets, there exists a constant $D>0$ such that $||\mathbf{f}(\mathbf{x}(t)) + \mathbf{G}(\mathbf{x}(t))\mathbf{u}(t)|| \leq D$, $\forall t\geq 0$. Hence, we now have a uniform lower bound for $||t_{k+1}-t_k||\geq \frac{\eta_h-\eta_l}{DL}$ for all $k=0,1,...$.} 
    Condition P1-1 is now proved.

    \textit{\underline{Proof of P1-2: }} If conditions C1-1 and C1-2 are satisfied, then all possible switching signals under Switching Strategy 1 satisfy P1-2. Hence, P1-2 is proved.
    
\end{proof}

\end{document}
